\documentclass[11pt,a4paper]{article}
\usepackage[utf8]{inputenc}
\usepackage[T1]{fontenc}
\usepackage{lmodern}
\usepackage[a4paper, margin=0.68in]{geometry}
\usepackage{amsmath, amssymb}
\usepackage{booktabs}
\usepackage{tabularx}
\usepackage{enumitem}
\usepackage{ragged2e}
\usepackage[colorlinks=true, linkcolor=black, urlcolor=blue, citecolor=blue]{hyperref}
\usepackage{titlesec}
\titlespacing*{\section}{0pt}{4pt}{2pt}
\titleformat{\section}{\normalfont\large\bfseries}{\thesection.}{0.5em}{}
\setlength{\parindent}{0pt}
\setlength{\parskip}{2.5pt}
\newcolumntype{Y}{>{\RaggedRight\arraybackslash}X}
\pagestyle{empty}

\begin{document}

\begin{center}
{\large\bfseries A Spiking State-Space Detector for Event-Based Micro-UAV Perception on Intel Loihi 2}\\[3pt]
\small Dr Will Midgley (PI) \& Benas Vaiciulis (student researcher) \,$\cdot$\, UNSW Sydney \,$\cdot$\,
PyTorch $\rightarrow$ NIR $\rightarrow$ Loihi 2\\[2pt]
\small\textbf{Request:} Neuromorphic Research Cloud (vLab) access now; hardware only after a positive
vLab benchmark.
\end{center}
\vspace{-2pt}
\hrule
\vspace{4pt}

\textbf{Objective.} Map a \emph{spiking state-space} object detector for event cameras onto \textbf{Intel Loihi 2},
targeting the size-weight-and-power (SWaP) budget of micro-UAVs. The spiking model is built and training. It is a
controlled swap into a \emph{completed and fully benchmarked} non-spiking reference pipeline, so the experiment
measures one thing cleanly: \textbf{what does spiking a state-space detector cost, where does that cost come
from, and what does neuromorphic hardware give back?}

%---------------------------------------------------------------
\section{The accuracy gap has closed; its cost has not been decomposed}
On Prophesee Gen1, the standard event-camera detection benchmark, spiking detectors have caught up: against
47.7 mAP for the non-spiking S5-RVT~[1], the hybrid HsVT~[2] reports up to 47.8 and the spiking SpikeDet~[3] up to
47.6. The neuromorphic case therefore rests on energy and latency, not accuracy. \emph{Where} the cost of
spiking comes from is unmeasured. SpikeYOLO~[4] and EMS-YOLO~[5] compare against an ANN of equivalent architecture, so the total cost is known for those designs; no published detector separates
it into the cost of the neuron dynamics, of sparsity and of binarisation, and none uses a selective state-space
temporal core.

\begin{table}[htbp]
\centering\small
\setlength{\tabcolsep}{4.5pt}
\begin{tabularx}{\textwidth}{@{}Y l c c@{}}
\toprule
\textbf{Model (Gen1 test, COCO mAP; ranges over model sizes)} & \textbf{Type} & \textbf{mAP} & \textbf{Params} \\
\midrule
S5-RVT, Zubi\'c \emph{et al.}, CVPR 2024~[1] & ANN (SSM) & 47.7 & 18.2\,M \\
\textbf{EventSSM} (ours) --- ResNet-18 conv $+$ Mamba temporal & ANN (SSM) & 46.2 & 19.2\,M \\
\textbf{PureSSM} (ours) --- BiMamba spatial $+$ Mamba temporal & ANN (SSM) & 46.4 & \textbf{16.3\,M} \\
\midrule
HsVT, ICML 2025~[2] --- spiking temporal block & hybrid SNN & 44.9--47.8 & --- \\
SpikeDet, preprint v5 2026 (v1: SpikSSD, 40.8)~[3] & SNN & 46.5--47.6 & --- \\
SpikeYOLO, ECCV 2024~[4] & SNN & 38.5--40.4 & 13.2 / 23.1\,M \\
EMS-YOLO, ICCV 2023~[5] & SNN & 26.7--31.0 & 6.2 / 14.4\,M \\
\midrule
\textbf{Proposed (this project)} & \textbf{hybrid SNN} & \textbf{training} & 16.3\,M \\
\bottomrule
\end{tabularx}
\end{table}

\vspace{-4pt}
\textbf{Scope, stated plainly.} We spike the \emph{temporal} readout of an already-benchmarked non-spiking
detector (PureSSM) and change nothing else. The result is a \textbf{hybrid} (spiking temporal path, non-spiking
spatial backbone), as in HsVT~[2], and \emph{not} a like-for-like competitor to the fully spiking
detectors above. Three readouts share one leaky integrate-and-fire membrane --- \emph{analog} (dense),
\emph{graded} (sparse, value-carrying) and \emph{binary} --- so the steps PureSSM $\rightarrow$ analog
$\rightarrow$ graded $\rightarrow$ binary isolate the cost of the \textbf{neuron dynamics}, of \textbf{sparsity}
and of \textbf{binarisation}. A state-space neuron already supplies recurrence, so we expect binarisation to cost
most --- the cost Loihi 2's \textbf{graded spikes} target. 25k-step pilots agree (validation AP, compressed
schedule, not comparable to the table): PureSSM 35.1, analog 35.1, graded 34.0, binary 30.5. The equal-budget
100k-step ladder is training now.

%---------------------------------------------------------------
\section{What is already built}
We reproduced S5-RVT exactly (Gen1 test AP $=47.7$) and built two detectors as \emph{controlled} backbone swaps
into that pipeline (same YOLO-PAFPN neck, YOLOX head, loss, data and evaluator). On a calibrated, idle-GPU-guarded
harness PureSSM is the \textbf{smallest} (16.3\,M), \textbf{cheapest} (10.0\,GFLOP), \textbf{most accurate per
FLOP} (4.63 mAP/GFLOP) and \textbf{most rate-robust} model: it retains \textbf{69.7\,\%} of its accuracy under a
true $10\times$ event-rate change (S5-RVT: 62.2\,\%).

\textbf{The spiking model is built.} A leaky integrate-and-fire readout, trained with a surrogate gradient, wraps
each of PureSSM's Mamba temporal blocks unmodified. With no spiking stage it is numerically identical to PureSSM;
streaming and full-sequence inference differ in under 0.1\,\% of elements; a 428-test suite covers the cell, the
integration and the cross-clip state; and pilots of all three readouts passed our pre-registered
stability criteria (finite loss, no silent or saturated layer, rising accuracy).

%---------------------------------------------------------------
\section{Why Loihi 2 specifically}
\begin{enumerate}[leftmargin=1.4em, itemsep=1.5pt, topsep=1.5pt]
  \item \textbf{The temporal core is a diagonal linear recurrence} $x_k=\bar{A}_k x_{k-1}+\bar{B}_k u_k,\;y_k=C_k x_k$
  --- a bank of independent scalar states. Loihi 2's \textbf{programmable neuron microcode}~[7] implements this class
  of dynamics directly, instead of a fixed LIF primitive. Meyer \emph{et al.}~[6] mapped a diagonal
  SSM (S4D) to Loihi 2 and, for token-by-token processing --- the regime of a streaming event camera --- measured
  ${\sim}1000\times$ lower energy and ${\sim}75\times$ lower latency than a recurrent S4D on a Jetson Orin Nano.
  \item \textbf{Event sparsity is wasted on a GPU but rewarded on Loihi 2.} We measured that sparse event input
  yields \emph{no} GPU speed-up, as dense kernels process idle pixels regardless; on an event-driven substrate idle
  pixels cost no energy. The 1\,Mpx event-detection paper names neuromorphic hardware as the route to the
  sparsity its own dense pipeline wastes~[8].
  \item \textbf{Graded spikes address the binary bottleneck.} Integer-valued spiking detectors~[3,\,4] train with
  multi-bit activations and expand each into several binary spikes at inference. Loihi 2's graded spikes carry an
  integer payload natively~[7] --- the hardware counterpart of our graded readout --- so that accuracy recovery is
  a Loihi 2 capability, not a software workaround.
\end{enumerate}

With PureSSM's rate robustness, this leads from ${\sim}0.65$\,J/frame on a GPU toward UAV-scale power.

%---------------------------------------------------------------
\section{Approach and milestones}
The trained model keeps Mamba's \emph{selective} step: the decay $\bar{A}_k=\exp(\Delta_k A)$ is recomputed from
the input at every step, whereas the S4D port~[6] relies on fixed dynamics. The mapping therefore needs on-chip
computation of the decay or a \textbf{fixed-step} (S4D/S5-style) variant retrained through the same pipeline; our
rate study indicates a fixed step need not cost robustness (at a true $10\times$ rate the fixed-step S5-RVT retains
62.2\,\%, our selective EventSSM 63.0\,\%). The non-spiking spatial backbone would run on a conventional
accelerator beside the chip. Following Intel's guidance to design locally first, we train in PyTorch
(surrogate gradients), export through \textbf{NIR}, the vendor-neutral neuromorphic intermediate
representation, and benchmark on vLab before any hardware request. \emph{As the Lava repositories were archived in May 2026, we treat NIR as the export
boundary so the model retargets cleanly to a successor SDK; guidance on Intel's preferred path is welcome.}

\vspace{-1pt}
\begin{center}\small
\begin{tabularx}{\textwidth}{@{}l Y@{}}
\toprule
\textbf{M1} (wk 1--4) & Spiking SSM cell (SSM $+$ LIF readout $+$ surrogate gradient), unit-tested. \textbf{Done.} \\
\textbf{M2} (wk 5--10) & Integrated into the frozen detector; trained on Gen1; accuracy, firing rate and sparsity.
\textbf{Integrated; 100k-step ladder training.} \\
\textbf{M3} (wk 11--16) & Fixed-step or on-chip-decay temporal block; NIR export $\rightarrow$ Loihi runtime
(Lava or successor); \textbf{vLab benchmark} --- accuracy, energy (SOPs), latency vs GPU/embedded baselines. \\
\textbf{M4} (gated on M3) & On-chip validation on Loihi 2 (Kapoho Point); energy--accuracy--latency
characterisation; publication. \\
\bottomrule
\end{tabularx}
\end{center}

\vspace{2pt}
\textbf{Readiness against the INRC checklist --- stated plainly.}
(i) \emph{Project plan} --- this document.
(ii) \emph{Built and tested an SNN} --- \textbf{partly met}: built, tested and training in simulation (Section~2),
since ANN and RNN experience alone does not prepare a team for Loihi.
(iii) \emph{Loihi toolchain} --- not yet; M3, via NIR.
(iv) \emph{Evaluated on vLab} --- M3, the gate on any hardware request. \emph{Hence vLab access now:} the cloud is
where we would port and validate the SNN; hardware follows only once M3 justifies it.

\vspace{2pt}
\textbf{Contributed back.} An open-source, NIR-exportable spiking state-space detector --- to our knowledge the
first for event cameras on Loihi 2; the first decomposition of the cost of spiking a detector into neuron
dynamics, sparsity and binarisation; and a Gen1 benchmark point for INRC event-vision work.

%---------------------------------------------------------------
\vspace{2pt}
{\footnotesize
\textbf{References.}
[1] N. Zubi\'c, D. Gehrig, M. Gehrig, and D. Scaramuzza, ``State space models for event cameras,'' in
\emph{Proc. CVPR}, 2024, pp.~5819--5828.\hspace{0.5em}
[2] Q. Xu, J. Deng, J. Shen, B. Chen, H. Tang, and G. Pan, ``Hybrid spiking vision transformer for object
detection with event cameras,'' in \emph{Proc. ICML}, 2025.\hspace{0.5em}
[3] Y. Fan, C. Liu, M. Li, D. Liu, Y. Su, Y. Liu, and W. Zhang, ``SpikeDet: Better firing patterns for accurate
and energy-efficient object detection with spiking neural networks,'' arXiv:2501.15151, v5, 2026.\hspace{0.5em}
[4] X. Luo, M. Yao, Y. Chou, B. Xu, and G. Li, ``Integer-valued training and spike-driven inference spiking
neural network for high-performance and energy-efficient object detection,'' in \emph{Proc. ECCV},
2024.\hspace{0.5em}
[5] Q. Su, Y. Chou, Y. Hu, J. Li, S. Mei, Z. Zhang, and G. Li, ``Deep directly-trained spiking neural networks
for object detection,'' in \emph{Proc. ICCV}, 2023.\hspace{0.5em}
[6] S.\,M. Meyer, P. Weidel, P. Plank, L. Campos-Macias, S.\,B. Shrestha, P. Stratmann, and M. Richter, ``A
diagonal structured state space model on Loihi 2 for efficient streaming sequence processing,''
arXiv:2409.15022, 2024.\hspace{0.5em}
[7] G. Orchard, E.\,P. Frady, D. Ben Dayan Rubin, S. Sanborn, S.\,B. Shrestha, F.\,T. Sommer, and M. Davies,
``Efficient neuromorphic signal processing with Loihi 2,'' in \emph{Proc. IEEE SiPS}, 2021.\hspace{0.5em}
[8] E. Perot, P. de Tournemire, D. Nitti, J. Masci, and A. Sironi, ``Learning to detect objects with a
1 megapixel event camera,'' in \emph{Proc. NeurIPS}, 2020.
\par}

\end{document}
